\documentclass[../../main.tex]{subfiles}
\begin{document}

\begin{center}
  {\fontsize{48pt}{58pt}\selectfont \textbf{2008}\normalsize\textbf{年度}}
\end{center}

\begin{table}[htb]
  \centering
  \caption{出題テーマと難易度}
  \begin{tabular}{|c|p{9cm}|c|}
    \hline
    \textbf{問題番号} & \textbf{テーマ} & \textbf{難易度} \\
    \hline
    第1問 & 一次変換で移る直線と領域 & \\
    \hline
    第2問 & カードの色を変える操作と確率 & \\
    \hline
    第3問 & 正八面体の投影図と回転体の体積 & \\
    \hline
    第4問 & 放物線上の二点を結ぶ線分と中点の高さ & \\
    \hline
    第5問 & すべての桁が1の整数と約数 & \\
    \hline
    第6問 & 媒介変数で表された曲線と面積 & \\
    \hline
  \end{tabular}
\end{table}

\begin{table}[htb]
  \centering
  \caption{解答}
  \renewcommand{\arraystretch}{1.4}
  \begin{tabular}{|c|c|p{9cm}|}
    \hline
    \textbf{問題番号} & \textbf{小問} & \textbf{答え} \\
    \hline
    第1問 & (1) & $a_{n+1}=b_n$，$b_{n+1}=-\dfrac{1}{2}a_n+\dfrac{3}{2}b_n$ \\
    \cline{2-3}
    & (2) & 領域を図示（解答参照） \\
    \hline
    第2問 & (1) & $n$が奇数なら$0$，偶数なら$\dfrac{1}{4}\left(\dfrac{3}{4}\right)^{\frac{n}{2}-1}$ \\
    \cline{2-3}
    & (2) & $n=1$または偶数なら$0$，奇数$n\ge3$なら
    $\dfrac{1}{18}\left(\dfrac{17}{18}\right)^{\frac{n-3}{2}}$ \\
    \hline
    第3問 & (1) & 投影図を図示（解答参照） \\
    \cline{2-3}
    & (2) & $\dfrac{5\sqrt{6}}{54}\pi$ \\
    \hline
    第4問 & (1) & $h=\dfrac{1}{4}\left(m^2+\dfrac{L^2}{1+m^2}\right)$ \\
    \cline{2-3}
    & (2) & 最小値は$0<L\le1$のとき$\dfrac{L^2}{4}$，
    $L\ge1$のとき$\dfrac{2L-1}{4}$ \\
    \hline
    第5問 & (1) & 1を$3^m$個並べた整数は$3^m$で割り切れ，
    $3^{m+1}$で割り切れない（証明問題） \\
    \cline{2-3}
    & (2) & 1を$n$個並べた整数が$27$で割り切れることと，
    $n$が$27$で割り切れることは同値（証明問題） \\
    \hline
    第6問 & -- & $\dfrac{32}{9}$ \\
    \hline
  \end{tabular}
\end{table}

\paragraph{試験形態}
6問．試験時間は要確認．

\end{document}
